\chapter{Удиртгал}
\label{ch:introduction}

%======================================================================
% ЭНЭ ФАЙЛД БАЙГАА БҮХ ЗҮЙЛ БАГШИД ҮЗҮҮЛЭХ НООРОГТ ХЭВЛЭГДЭНЭ (бүтэн ноорогт ч мөн).
% Зөвхөн бүтэн ноорогт орох хэсэг тусдаа файлд: Chapters/full/ch1_*.tex (\fullinput мөрүүд).
%
% 2026.11.16 (2), хэрэглэгчийн шүүмжээр дахин бичсэн: үндэслэлийг вэб сайтаар биш, сургуульд
% ирсэн хүн ба ажилтны бэрхшээлээр эхлүүлсэн; авсан мэдээллийг тоочихгүй (тоо 4.1-ийн хүснэгтэд);
% мэдээллийн эх сурвалж = вэб сайт ба сургуулийн албад; «мэдлэгийн мод» нь арга тул зорилт,
% ач холбогдлоос авсан (3.3-т тайлбарласан). Өмнөх хувилбар: thesis/archive-20261116/.
% Бүтэц, үг хэллэг: papers/28-mn-thesis-style-20261111.mn.md (Монгол 7 дипломын тулгалт).
%======================================================================

\section{Үндэслэл, ач холбогдол}
\label{sec:intro-background}

Шинэ Монгол Технологийн Коллежид ирсэн оюутан, элсэгч, эцэг эх, зочин хэрэгтэй мэдээллээ нэгдүгээр давхрын ресепшний ажилтан болон сургуулийн бусад ажилтнаас биечлэн асууж авдаг. Сургуулийн мэдээлэл нэг газарт төвлөрөөгүй. Нэг хэсэг нь сургуулийн вэб сайтад, нэг хэсэг нь сургуулийн дотоод баримт бичигт байдаг.

Хүмүүс юу асуудгийг ресепшний ажилтнаас тодруулав. Ажилтны мэдээлснээр асуулт хоёр бүлэгт хуваагдана: өдөр тутам ирдэг хүмүүсийн асуулт ба нээлттэй өдөрлөг, сургуулийн танилцуулгын үеэр ирдэг хүмүүсийн асуулт (\tref{tab:reception_topics}). Өдөр тутмын таван сэдвийн гурав нь хэнтэй холбогдох, хаана очихыг тодруулсан асуулт байна.

\begin{table}[ht]
\centering
\caption{Ресепшнд хамгийн их асуудаг асуултын сэдэв (ресепшний ажилтны мэдээлснээр, 2026.11.12)}
\label{tab:reception_topics}
\begin{tabular}{|p{6.8cm}|p{6.8cm}|}
\hline
\textbf{Өдөр тутам} & \textbf{Нээлттэй өдөрлөг, танилцуулгын үеэр} \\
\hline
Сургалтын төлбөр & Сургалтын орчин \\
\hline
Эцэг эх багштай уулзахаар утасны дугаарыг нь асуух & Төлбөрийн асуудал \\
\hline
Ажлын байр сонирхогч анкетаа хаана бөглөх & Харилцаа, хандлага \\
\hline
Ярилцлагад ирсэн хүн хэдэн давхарт, хаана очих & Дотуур байр байдаг эсэх \\
\hline
Гадны хүүхдэд давтлага өгдөг эсэх & \\
\hline
\end{tabular}
\end{table}

Энэ нь нэг ажилтны ажиглалтад тулгуурласан бөгөөд сэдэв тус бүрийн давтамжийг тоолоогүй.
% ӨӨРИЙН ҮГЭЭР: ажилтны албан тушаалыг нэм. Тоолох хуудсаар нэг долоо хоногийн тоо гарвал энд нэм.

Сургалтын албаны ажилтнуудтай хийсэн уулзалтаар хоёр бэрхшээл тодорхой болов. Нэгдүгээрт, ажилтнууд ижил мэдээллийг өөр өөр хүнд дахин дахин тайлбарладаг. Хоёрдугаарт, нэгдүгээр давхарт цаасан мэдээлэл байрлуулсан боловч хүмүүс түүнээс хариултаа олж чаддаггүй, ажилтнаас дэлгэрүүлж асуудаг. Ажилтнууд мэдээлэл өгөх систем нь хариулт олдоогүй үед хаанаас авахыг заадаг, ресепшнээс гадна гар утаснаас ч нээгддэг байвал зохино гэж зөвлөсөн.
% ӨӨРИЙН ҮГЭЭР: хэнтэй (албан тушаал) уулзсаныг нэм. Карьерын төв, олон нийтийн алба зэрэг
% бусад албатай уулзсан бол тэдний хэлснийг энд нэг, хоёр өгүүлбэрээр нэм (агент зохиохгүй).

Ийм бэрхшээл бусад сургуульд ч ажиглагдсан. Тухайлбал, Милано-Бикокка их сургуулийн зургаан оюутантай хийсэн ярилцлагад оюутнууд анги танхим, хичээлийн хуваарь, шалгалтын өрөөний мэдээллийг олоход хэцүү гэж хариулсан \cite{antico2024unimib}. Burch олон улсын их сургуулийн оюутнууд мэдээллээ олж чадаагүй үедээ энгийн асуултаар захиргааны ажилтанд ханддаг бөгөөд энэ нь ажилтны ачааллыг нэмдэг гэж зохиогчид тэмдэглэсэн \cite{matic2026burch}. Хоёр ажлын зохиогчид асуултад сургуулийн баримтаас хайж хариулдаг туслах систем хийж, энэ асуудлыг шийдэхийг зорьсон.

Ийм туслахыг Монгол хэл дээр хийхэд хоёр бэрхшээл тулгардаг. Нэгдүгээрт, Монгол хэл нь яриа таних загвар сургах өгөгдөл багатай хэлний тоонд ордог. Олон хэлний Whisper large-v2 загвар Монгол яриаг танихад үгийн алдааны хувь (WER, \ref{sec:lit-asr}-р хэсэг) 100\%-иас давсан \cite{radford2023whisper}. Common Voice 17 өгөгдөлд Халх Монгол хэлний сургалтын яриа ердөө 3 цаг бий \cite{mengke2025khalkha}. Хоёрдугаарт, сургуулийн нэрийн өмнөөс өгөх хариулт үнэн зөв байх ёстой. Хэлний загвар бодит бус мэдээлэл зохиох тохиолдол гардаг тул хариултыг баримтын сангаас хайж, олдсон эх сурвалжид тулгуурлан өгөх шаардлагатай \cite{lewis2020rag}.

Иймд сургуулийн тархай бутархай мэдээллийг нэг бүтцэд оруулж, Монгол хэлээр асуусан асуултад түүнээс хариулдаг, хариулт байхгүй үед хэнд, хаана хандахыг зааж өгдөг туслах систем боловсруулах шаардлага тодорхой болов.

Энэ ажлын практик ач холбогдол нь сургуулийн ажилтнуудын ижил мэдээллийг давтан тайлбарлах ачааллыг бууруулах, сургуульд ирсэн хүн хэрэгтэй мэдээллээ нэг газраас авах боломжийг бүрдүүлэхэд оршино. Судалгааны ач холбогдол нь сургуулийн албан мэдээллийг мэдлэгийн мод хэлбэрээр зохион байгуулж, түүнд суурилсан хайлтаар хариулдаг системийг Монгол хэл дээр хэрэгжүүлэн, ресепшнд асуудаг бодит асуултаар үнэлэхэд оршино.

\section{Зорилго, зорилт}
\label{sec:intro-goals}

\subsection{Судалгааны зорилго}

Энэ судалгааны зорилго нь сургуулийн ресепшнд ирсэн хүний Монгол хэлээр ярьж эсвэл бичиж асуусан асуултад \textbf{сургуулийн албан мэдээлэлд тулгуурлан хариулдаг туслах системийг боловсруулахад} оршино.

\subsection{Судалгааны зорилтууд}

Энэ зорилгод хүрэхийн тулд дөрвөн зорилтыг дэвшүүлэв:
\begin{enumerate}
    \item Ресепшнд ирдэг хүмүүсийн мэдээллийн хэрэгцээ, системд тавигдах шаардлагыг сургуулийн ажилтнуудтай уулзаж тодорхойлох.
    \item Сургуулийн албан мэдээллийг вэб сайтаас болон сургуулиас авч, мэдлэгийн мод хэлбэрээр зохион байгуулах.
    \item Дуугаар болон бичгээр асуусан асуултад эх сурвалжтай хариулт өгч, хариулт олдоогүй үед хандах алба, өрөөг заадаг системийг боловсруулж хэрэгжүүлэх.
    \item Системийг ресепшнд асуудаг бодит асуултаар үнэлэх.
\end{enumerate}
% Зорилт 1: карьерын төв, ресепшний ажилтантай уулзсаны дараа уулзсан албадыг нэрээр нь бич.
% 2026.11.16 (3), хэрэглэгч: мэдлэгийн мод бол RAG-ийн суурь — 2-р зорилт, ач холбогдолд буцааж оруулсан;
% үнэлгээ тусдаа 4-р зорилт хэвээр (үнэлгээг «мэдлэгийн модонд тулгуурласан» гэж нэрлэхгүй).

\fullinput{Chapters/full/ch1_research_questions}

\section{Судалгааны хамрах хүрээ}
\label{sec:intro-scope}

Систем зөвхөн сургуулийн албан мэдээллээр хариулна. Мэдээллийг сургуулийн вэб сайтаас болон сургуулийн зүгээс авсан. Асуулт, хариулт Монгол хэлээр байна. Дүн, ирц зэрэг хувь хүний мэдээлэл болон ажилтны хувийн утасны дугаар, хувийн имэйл хаягийг энэ ажилд хамруулаагүй.

% 2026.11.16 (багшийн зөвлөгөө): хэрэглээний тохиолдлын диаграммыг 3.2-оос энд, эхний бүлэгт шилжүүлсэн.
Системийг хэн, ямар зорилгоор ашиглахыг хэрэглээний тохиолдлын диаграммаар \fref{fig:usecase}-д үзүүлэв. Диаграммын зүүн талд системийг ашиглах зочдыг, баруун талд гадны хоёр үйлчилгээ болон сургуулийн ажилтныг байрлуулсан. Сургуульд ирсэн зочин асуултаа бичгээр эсвэл дуугаар тавьж, эх сурвалж, огноо бүхий хариулт авна. Хариулт олдоогүй тохиолдолд хэнд, хаана хандахыг мэднэ. Яриаг текст болгох, текстийг дуу болгох үйлдлийг гадны үйлчилгээ гүйцэтгэнэ. Сургуулийн ажилтан мэдээллийг шинэчлэн нийтлэх үйлдлийг цаашид идэвхжүүлэхээр төлөвлөсөн.

\begin{figure}[ht]
\centering
\adjustbox{max width=\textwidth}{% Хэрэглээний тохиолдлын (use case) диаграмм (TikZ).
% Эх: demo_server.py-ийн нээлттэй API (/api/route, /api/voice/transcribe, /api/voice/speak,
% /api/admin/*) ба tools/import_school_tables.py. Оролцогчдыг 2026.11.12-ны ресепшний
% мэдээллээр хоёр бүлэгт хуваасан: өдөр тутмын зочин, нээлттэй өдөрлөгийн зочин.
% Байрлалыг (x,y) координатаар өгсөн.
\begin{tikzpicture}[
    font=\footnotesize,
    uc/.style={draw, ellipse, align=center, text width=3.3cm, minimum height=0.95cm, inner sep=1pt},
    actor/.style={align=center, font=\footnotesize},
    extbox/.style={draw, align=center, text width=2.5cm, minimum height=0.9cm, fill=black!8},
    line/.style={semithick},
    dep/.style={dashed, -{Stealth[length=2mm]}},
    gen/.style={-{Triangle[open, length=2.6mm, width=2.6mm]}, semithick},
    lab/.style={font=\scriptsize, inner sep=1.5pt, fill=white}]

% --- хүний дүрс зурах макро: #1 нэр, #2 (x,y), #3 бичиг
\newcommand{\stickman}[3]{%
  \begin{scope}[shift={(#2)}]
    \draw[semithick] (0,0.55) circle (0.17);
    \draw[semithick] (0,0.38) -- (0,-0.1) (-0.3,0.2) -- (0.3,0.2) (0,-0.1) -- (-0.25,-0.5) (0,-0.1) -- (0.25,-0.5);
    \node[actor, anchor=north] at (0,-0.55) {#3};
    \coordinate (#1) at (0.35,0.1);
    \coordinate (#1-l) at (-0.35,0.1);
  \end{scope}}

% --- системийн хүрээ
\draw[semithick] (2.4,2.1) rectangle (11.0,-9.6);
\node[font=\footnotesize\bfseries] at (6.7,1.75) {Сургуулийн ухаалаг туслах};

% --- хэрэглээний тохиолдлууд
\node[uc] (typed)  at (4.6,0.7)   {Бичгээр асуух};
\node[uc] (voice)  at (4.6,-0.8)  {Дуугаар асуух};
\node[uc] (answer) at (4.6,-2.5)  {Эх сурвалж, огноотой\\хариулт авах};
\node[uc] (refer)  at (4.6,-4.2)  {Хэнд, хаана хандахыг\\мэдэх};
\node[uc] (listen) at (4.6,-5.9)  {Хариултыг дуугаар\\сонсох};
\node[uc] (stt)    at (8.9,-0.8)  {Яриаг текст болгох};
\node[uc] (tts)    at (8.9,-5.9)  {Текстийг дуу болгох};
\node[uc] (import) at (8.9,-7.4)  {Сургуулийн хүснэгт\\оруулах};
\node[uc] (publish) at (8.9,-8.8) {Мэдээллийн шинэ хувилбар\\шалгаж нийтлэх};

% --- оролцогчид
\stickman{guest}{0.4,-2.4}{Зочин}
\stickman{daily}{-3.2,-0.4}{Өдөр тутмын зочин\\{\scriptsize оюутан, эцэг эх,}\\{\scriptsize ажил горилогч}}
\stickman{open}{-3.2,-4.4}{Нээлттэй өдөрлөгийн\\зочин\\{\scriptsize элсэгч, эцэг эх}}
\stickman{staff}{13.4,-8.0}{Сургуулийн\\ажилтан}
\node[extbox] (gemini) at (13.4,-0.8) {Яриа таних\\үйлчилгээ\\{\scriptsize Gemini 3.5 Transcribe}};
\node[extbox] (edge)   at (13.4,-5.9) {Дуу үүсгэх\\үйлчилгээ\\{\scriptsize Edge TTS}};

% --- холбоос
\foreach \u in {typed, voice, answer, refer, listen} \draw[line] (guest) -- (\u.west);
\draw[gen] (daily) -- (guest-l);
\draw[gen] (open) -- (guest-l);
\draw[dep] (voice) -- node[lab] {«include»} (stt);
\draw[dep] (listen) -- node[lab] {«include»} (tts);
\draw[dep] (refer) -- node[lab, right] {«extend»} (answer);
\draw[line] (stt) -- (gemini);
\draw[line] (tts) -- (edge);
\draw[line] (staff-l) -- (import.east);
\draw[line] (staff-l) -- (publish.east);
\end{tikzpicture}
}
\caption{Системийн хэрэглээний тохиолдлын диаграмм}
\label{fig:usecase}
\end{figure}

Багш, ажилтны нэр, ажлын өрөө, сургуулийн имэйл хаяг, хичээлийн хуваарийг сургуулийн өгсөн хэмжээгээр ресепшний дэлгэц дээр харуулна. Системийг сургуулийн гаднаас гар утсаар нээдэг болгох үед эдгээрийг харуулах эсэхийг сургуультай тохирч шийднэ.

\fullinput{Chapters/full/ch1_structure}
